% GENERATED FILE. Edit canonical lessons, then run npm run generate:books.
% canonical-source-hash: fnv1a64:0eb83a7d36e770cc
% canonical-lessons: RU-C362-povyshenie, RU-C362-premiya, RU-C362-navyk, RU-C362-specialist, RU-C362-uvolnyat

\chapter{Getting a promotion}
\label{ch:ru-getting-a-promotion}
\hlchaptermodality{362}

\begin{chapteropening}
\textbf{By the end of this chapter:} \emph{I can talk about getting on at work: a promotion, a bonus, my skills, being a specialist, and being dismissed.}

Complete the last lesson of chapter 362: I can talk about getting on at work: a promotion, a bonus, my skills, being a specialist, and being dismissed.
\end{chapteropening}

\section[\texorpdfstring{\ru{повышение} (povyshéniye)}{повышение (povyshéniye)}]{\ru{повышение} (povyshéniye) — a promotion; a rise}

\label{lesson:RU-C362-povyshenie}



\begin{quote}
Before the new one: say the Russian for income, then the Russian for a deal.
\end{quote}

\subsection*{You'll want to know: \ru{повышение}}
\textbf{\ru{повышение}} (\emph{povyshéniye}) — \textquotedblleft{}a promotion; a rise\textquotedblright{}.

\begin{quote}
\ru{Он} \ru{получил} \ru{повышение} \ru{на} \ru{работе}.
\end{quote}

\begin{quote}
\emph{On poluchíl povyshéniye na rabóte.}
\end{quote}

\begin{quote}
He got a promotion at work.
\end{quote}

Neuter. \emph{\ru{Повышение} \ru{зарплаты}} — a pay rise. From \emph{\ru{повышать}}.

\begin{cousinweb}
\textbf{\ru{повышение}} is built on \emph{\ru{повысить}} \textquotedblleft{}to raise\textquotedblright{}, from \emph{\ru{выше}} \textquotedblleft{}higher\textquotedblright{}, the comparative of \emph{\ru{высокий}}, inherited from Proto-Slavic \emph{*vysokŭ} \textquotedblleft{}high, tall\textquotedblright{} (EtymDB).
\end{cousinweb}

\begin{grammarlens}[title={What you've built}]
One more word for telling, explaining and describing at B1.
\end{grammarlens}

\subsection*{Your turn}
\begin{itemize}
\raggedright
  \item \emph{Say it:} \emph{povyshéniye}
  \item \emph{Say it:} \emph{povyshéniye}, once more
  \item \emph{Say it:} a sentence of your own with \emph{povyshéniye}
  \item \emph{Recall:} say the Russian for income, then the Russian for a deal, then say \emph{povyshéniye} again
\end{itemize}

\begin{tcolorbox}[breakable,colback=teal!4,colframe=teal!35!black,title={Before you move on}]
What does \ru{повышение} mean? (\textquotedblleft{}A promotion; a rise\textquotedblright{}.) Say it once more.
\end{tcolorbox}
\section[\texorpdfstring{\ru{премия} (prémiya)}{премия (prémiya)}]{\ru{премия} (prémiya) — a bonus; a prize}

\label{lesson:RU-C362-premiya}



\begin{quote}
Before the new one: say the Russian for a deal, then the Russian for a promotion.
\end{quote}

\subsection*{You'll want to know: \ru{премия}}
\textbf{\ru{премия}} (\emph{prémiya}) — \textquotedblleft{}a bonus; a prize\textquotedblright{}.

\begin{quote}
\ru{В} \ru{конце} \ru{года} \ru{все} \ru{получили} \ru{премию}.
\end{quote}

\begin{quote}
\emph{V kontsé góda vse poluchíli prémiyu.}
\end{quote}

\begin{quote}
At the end of the year everyone got a bonus.
\end{quote}

Feminine. A bonus at work, or a prize: \emph{\ru{литературная} \ru{премия}}.

\begin{cousinweb}
\textbf{\ru{премия}} is usually traced, through German \emph{Prämie}, to Latin \emph{praemium} \textquotedblleft{}a reward, a prize\textquotedblright{} (Lewis \& Short), built from \emph{prae} \textquotedblleft{}before\textquotedblright{} and \emph{emere} \textquotedblleft{}to take, to buy\textquotedblright{}.
\end{cousinweb}

\begin{grammarlens}[title={What you've built}]
One more word for telling, explaining and describing at B1.
\end{grammarlens}

\subsection*{Your turn}
\begin{itemize}
\raggedright
  \item \emph{Say it:} \emph{prémiya}
  \item \emph{Say it:} \emph{prémiya}, once more
  \item \emph{Say it:} a sentence of your own with \emph{prémiya}
  \item \emph{Recall:} say the Russian for a deal, then the Russian for a promotion, then say \emph{prémiya} again
\end{itemize}

\begin{tcolorbox}[breakable,colback=teal!4,colframe=teal!35!black,title={Before you move on}]
What does \ru{премия} mean? (\textquotedblleft{}A bonus; a prize\textquotedblright{}.) Say it once more.
\end{tcolorbox}
\section[\texorpdfstring{\ru{навык} (návyk)}{навык (návyk)}]{\ru{навык} (návyk) — a skill}

\label{lesson:RU-C362-navyk}



\begin{quote}
Before the new one: say the Russian for a promotion, then the Russian for a bonus.
\end{quote}

\subsection*{You'll want to know: \ru{навык}}
\textbf{\ru{навык}} (\emph{návyk}) — \textquotedblleft{}a skill\textquotedblright{}.

\begin{quote}
\ru{Для} \ru{этой} \ru{работы} \ru{нужны} \ru{компьютерные} \ru{навыки}.
\end{quote}

\begin{quote}
\emph{Dlya étoy rabóty nuzhný kompyúternyye návyki.}
\end{quote}

\begin{quote}
This job needs computer skills.
\end{quote}

Masculine; usually plural, \emph{\ru{навыки}}.

\begin{cousinweb}
\textbf{\ru{навык}} is built from \emph{\ru{на}-} and the root of \emph{\ru{привыкать}} \textquotedblleft{}to get used to\textquotedblright{} and \emph{\ru{отвыкать}} \textquotedblleft{}to lose a habit\textquotedblright{}: a skill is something you have grown used to doing.
\end{cousinweb}

\begin{grammarlens}[title={What you've built}]
One more word for telling, explaining and describing at B1.
\end{grammarlens}

\subsection*{Your turn}
\begin{itemize}
\raggedright
  \item \emph{Say it:} \emph{návyk}
  \item \emph{Say it:} \emph{návyk}, once more
  \item \emph{Say it:} a sentence of your own with \emph{návyk}
  \item \emph{Recall:} say the Russian for a promotion, then the Russian for a bonus, then say \emph{návyk} again
\end{itemize}

\begin{tcolorbox}[breakable,colback=teal!4,colframe=teal!35!black,title={Before you move on}]
What does \ru{навык} mean? (\textquotedblleft{}A skill\textquotedblright{}.) Say it once more.
\end{tcolorbox}
\section[\texorpdfstring{\ru{специалист} (spetsialíst)}{специалист (spetsialíst)}]{\ru{специалист} (spetsialíst) — a specialist, an expert}

\label{lesson:RU-C362-specialist}



\begin{quote}
Before the new one: say the Russian for a bonus, then the Russian for a skill.
\end{quote}

\subsection*{You'll want to know: \ru{специалист}}
\textbf{\ru{специалист}} (\emph{spetsialíst}) — \textquotedblleft{}a specialist, an expert\textquotedblright{}.

\begin{quote}
\ru{Он} \ru{хороший} \ru{специалист} \ru{по} \ru{компьютерам}.
\end{quote}

\begin{quote}
\emph{On khoróshiy spetsialíst po kompyúteram.}
\end{quote}

\begin{quote}
He is a good computer specialist.
\end{quote}

Masculine, also of a woman. \emph{\ru{Специалист} \ru{по}} + dative — an expert in.

\begin{cousinweb}
\textbf{\ru{специалист}} is usually traced to German \emph{Spezialist}, built on Latin \emph{speciālis} \textquotedblleft{}particular\textquotedblright{}, from \emph{speciēs} \textquotedblleft{}a look, a kind\textquotedblright{} (Lewis \& Short).
\end{cousinweb}

\begin{grammarlens}[title={What you've built}]
One more word for telling, explaining and describing at B1.
\end{grammarlens}

\subsection*{Your turn}
\begin{itemize}
\raggedright
  \item \emph{Say it:} \emph{spetsialíst}
  \item \emph{Say it:} \emph{spetsialíst}, once more
  \item \emph{Say it:} a sentence of your own with \emph{spetsialíst}
  \item \emph{Recall:} say the Russian for a bonus, then the Russian for a skill, then say \emph{spetsialíst} again
\end{itemize}

\begin{tcolorbox}[breakable,colback=teal!4,colframe=teal!35!black,title={Before you move on}]
What does \ru{специалист} mean? (\textquotedblleft{}A specialist, an expert\textquotedblright{}.) Say it once more.
\end{tcolorbox}
\section[\texorpdfstring{\ru{увольнять} (uvolnyát)}{увольнять (uvolnyát)}]{\ru{увольнять} (uvolnyát) — to dismiss, to fire}

\label{lesson:RU-C362-uvolnyat}



\begin{quote}
Before the new one: say the Russian for a skill, then the Russian for a specialist.
\end{quote}

\subsection*{You'll want to know: \ru{увольнять}}
\textbf{\ru{увольнять}} (\emph{uvolnyát}) — \textquotedblleft{}to dismiss, to fire\textquotedblright{}.

\begin{quote}
\ru{Компания} \ru{увольняет} \ru{сто} \ru{человек}.
\end{quote}

\begin{quote}
\emph{Kompániya uvolnyáyet sto chelovék.}
\end{quote}

\begin{quote}
The company is laying off a hundred people.
\end{quote}

Imperfective; perfective \emph{\ru{уволить}}. \emph{\ru{Увольняться}} — to resign, to quit.

\begin{cousinweb}
\textbf{\ru{увольнять}} is formed from \emph{\ru{уволить}} with the suffix \emph{-\ru{ять}} (EtymDB). \emph{\ru{Уволить}} is built on \emph{\ru{воля}} \textquotedblleft{}will, freedom\textquotedblright{}, inherited from Proto-Slavic \emph{*voľa} (EtymDB): to dismiss someone was first to set them free from service.
\end{cousinweb}

\begin{grammarlens}[title={What you've built}]
One more word for telling, explaining and describing at B1.
\end{grammarlens}

\subsection*{Your turn}
\begin{itemize}
\raggedright
  \item \emph{Say it:} \emph{uvolnyát}
  \item \emph{Say it:} \emph{uvolnyát}, once more
  \item \emph{Say it:} a sentence of your own with \emph{uvolnyát}
  \item \emph{Recall:} say the Russian for a skill, then the Russian for a specialist, then say \emph{uvolnyát} again
\end{itemize}

\begin{tcolorbox}[breakable,colback=teal!4,colframe=teal!35!black,title={Before you move on}]
What does \ru{увольнять} mean? (\textquotedblleft{}To dismiss, to fire\textquotedblright{}.) Say it once more.
\end{tcolorbox}
